(1) 自旋极化意味着每个空间轨道只可占一粒子，基态占据 \(n=0,1,\dots,N-1\)。总能量 \(E_0=\hbar\omega\sum_{n=0}^{N-1}(n+1/2)=\boxed{\hbar\omega N^2/2}\)，与全部玻色子在基态时的 \(N\hbar\omega/2\) 明显不同。

(2) 最高占据单粒子能量 \(\varepsilon_F=\hbar\omega(N-1/2)\)。谐振态 \(\langle x^2\rangle_n=\hbar(n+1/2)/(m\omega)\)，故云团平均 \(\langle x^2\rangle=\hbar N/(2m\omega)\)。谐振维里定理给 \(E_{\rm kin}=E_{\rm pot}=E_0/2=\hbar\omega N^2/4\)。

(3) 半经典局域Fermi动量满足 \(p_F(x)=\sqrt{2m[\mu-m\omega^2x^2/2]}\)，在 \(|x|<x_F=\sqrt{2\mu/(m\omega^2)}\) 内的粒子密度为 \(\boxed{n(x)=p_F(x)/(\pi\hbar)}\)，外侧为0。积分 \(N=\int n(x)dx=\mu/(\hbar\omega)\)，所以对 \(N\gg1\) 取 \(\mu\simeq N\hbar\omega\)，和离散填充结果的主阶相同。

(4) 若慢速改变频率且维持绝热占据，\(E_0\propto\omega\)，而云团平方半径 \(\langle x^2\rangle\propto\omega^{-1}\)。即使所有粒子都倾向阱心，Pauli不相容使它们占据越来越高的单粒子激发态，零温仍拥有正的Fermi动能和有限云团宽度。
